购买系统的整体架构包含五条彼此关联的路径：调用方提交购买，所属分区保存业务决定，支付服务取得外部资金结果，消息服务传播已付款事实，履约服务接受并完成发货。它们共同对外呈现一次购买的进度，各自保存的事实却有明确归属。Alice 在页面上看到的订单，可能同时涉及购买分区里的 \code{C:b1:1}、支付平台里的 X1、消息日志里的两次 \code{EV:b1:1}，以及履约库里的 F501。系统设计首先要说明这些对象怎样连接，哪个保存域能够决定它们的存在。

图~\ref{fig:ch16-architecture} 把系统分成接入、业务权威、异步协作和外部效果四层。调用方先保存完整身份 $S$ 和原输入，入口完成认证与结构检查，再用路由缓存找到所属逻辑分区。业务组的确定性状态机在复制命令中改变会话、受理、付款责任、原结果和事件。支付工作者、发布器与发货扫描器从各自保存域枚举尚未完成的责任，沿原身份继续。它们可以换进程，持久责任随对应保存域恢复。

\CSixteenFigure{ch16-01-architecture}{完整系统按接入、业务权威、异步协作和外部效果分层，控制与派生路径单独标示}{fig:ch16-architecture}

图中的绿色保存域拥有业务决定权：购买分区决定受理与结果，履约数据库决定初次履约，平台决定各自资金与发货效果。蓝色执行组件负责搬运输入和证据，暂停以后可由后继者接替。控制面保存分区归属、配置与迁移资格，通过虚线下发位置和资格信息。缓存与图像存储位于右侧，前者提供可重建观察，后者承载带完整清单的恢复材料。这些关系确定了后续各章深入节点、复制和存储时所处的位置。

这里先交代第三部分到本部分的部署变化。第十四、十五章把决定日志 J 与购买库 W 分在两个恢复域：J 保存 $A_J/R_J$、会话和退款决定，W 按决定物化 P/C/T/O/EV，因而能够在 W 回退后由 J 重建。本部分选择另一条部署主线，将这些权威决定、业务关联和未完成责任共同放入每个逻辑分区的复制状态机。一次受理或结果命令在同一应用边界中改变相关记录，付款、发布和恢复工作者读取这份共同状态。$A_J/R_J$ 继续表示原受理与原结果，记号中的 J 指出决定的角色，本部分没有另外调用一套 J 服务来完成每次购买。高频页面查询从这份权威派生投影，再由 Redis 缓存。

第十四章的逻辑分片 $g$ 与这里的逻辑分区 b 指同一层的业务归属单位。例如，前面的 g7 可以在部署目录中登记为这里的 b1，导入时保存这一对应关系及全部原编号。G1、G2、G3 则是当前承载逻辑分区的副本组。更换 G1 的副本或把 b1 搬到 G2，都只改变承载与资格，已有 $S$、C 和平台键随原记录交接。

后文新受理采用 \code{C:b1:1} 一类编号，其前提是多数复制和在线迁移完整保留已提交的分配历史，分区序号与受理绑定共同提交。整组灾难回退可能同时丢掉序号和旧绑定，此时先执行第十四章的冻结与恢复准入：从跨故障域来源找回原 $S$、原 C 和原平台键，找不回的旧范围按持久关闭规则拒绝提交。核准开放真正的新意图后，才用不可回退来源授予的灾难代次区分新的编号空间，新编号共同包含稳定分区、灾难代次与该代次内序号，旧记录保留原字面。灾难代次防止新编号复用，恢复旧绑定的依据来自已保存的原历史；换一个代次不会重建已经丢失的 $S\to C$ 关系。

\subsection{功能边界和业务不变量}

本章使用一条独立的假设购买：Alice 的账户编号为 1，平台初始余额为 10000 分，商品 \code{paper-001} 的报价 Q7 为 CNY 3000 分，购买数量为 1。调用方用 \code{purchase.create.external.v2} 表示外部付款购买动作，在会话 E1 中保存客户端键 K1。完整身份继续采用
\[
S=(\texttt{account\_id},\texttt{action},\texttt{session\_id},\texttt{idempotency\_key})\text{。}
\]
同一 $S$ 的重试恢复同一购买，新一项购买由调用方确认新意图并保存新的身份。

系统允许付款先处于未知，再通过查询、回调或对账取得结论。付款成功以后可以先返回已付款、尚未接受履约，消息积压消退以后再推进到履约已接受。退款单独受理，每项退款有自己的 RF 身份，并关联原收款。组合商品还要先取得多个库存分区的联合预留决定，才允许真正发送付款。服务返回的每一种进度都对应一份已经保存的事实。

购买权威维持第一组不变量：每个完整 $S$ 绑定一份不可覆盖的原语义输入和一套实际 P/C/T。P 标识购买，C 标识付款责任，T 保存推动该责任完成的待办。组内换主、工作者更换、路由更新和分区迁移沿用原关联。原输入已经提交以后，后来同 $S$ 携带 3500 分的请求得到冲突，原 3000 分购买及其结果继续保存。

资金路径维持第二组不变量：同 C 的全部外部尝试取自保存的原平台范围、原平台键和原支付输入。未知结果继续占有原付款身份。可信成功在同一次权威提交中固定原平台结果、原付款结果、订单 O 和事件 EV。退款的成功金额与尚未排除成功的占用共同受原收款金额约束，未知退款也承担额度。

库存与履约分别维持自己的关系。一个库存分区的初始量等于可用量、准备占用和已提交预留之和，协调者的持久 commit 决定才授予组合购买的付款资格。履约库在同一事务中保存 Inbox 记录 I、业务守卫 B、实际履约编号 F 和发货意图。同一事件重投恢复原 F，不同事件若都声称接受同一订单的初次履约，还会在 B 上汇合。实际发货平台用原 dispatch 身份约束自己的外部效果。

这些安全关系可以在服务暂停时保持。进展则要求相关权威最终恢复提交能力，工作者得到调度机会，外部平台能终结并提供可核验结果，消息与履约域最终恢复处理。一个支付平台长期只提供 \code{pending}，系统就持续保存这份付款责任。面向用户的查询显示当前已知事实，运维侧则看到它的年龄、下一次查询安排与负责恢复者。

\subsection{负载和质量属性如何约束架构}

先给一组用于设计的负载假设。峰值入口每秒 5000 次，其中 2000 次为新 $S$，3000 次为已有 $S$ 的重试或查询。新 $S$ 需要仲裁首次受理和分配编号，同 $S$ 查询主要读取已存在的关系，付款恢复还可能访问平台。同样的入口请求数，在重试比例、未知支付比例和消息积压不同的情况下，会产生不同的后台压力。

\begin{table}[htbp]
\centering\small
\begin{tabularx}{\linewidth}{@{}p{.18\linewidth}p{.35\linewidth}X@{}}
\toprule
设计对象 & 本章采用的假设 & 对结构的要求\\
\midrule
峰值入口 & 每秒 5000 次，新 $S$ 占 40\% & 分开计量首次受理与已有操作读取，按分区限制新责任\\
受理响应 & 健康多数可达时，95\% 在 200 毫秒内取得受理或明确拒绝 & 受理不等待外部付款和发货，控制面不逐请求加锁\\
平台不可用 & 允许持续 10 分钟的未知与重试积压 & C/T 持久保存，恢复并发与新请求分别分配预算\\
消息恢复 & 20 分钟内排空一次 10 分钟积压 & 发布和消费分别留出恢复能力\\
部署故障 & 一个地域不可达时，其余两地域保存各组多数 & 业务组与控制组分别跨域放置，外部依赖单独评估\\
已提交事实 & 单副本退出后从完整多数事实继续 & 受理、原结果、会话、退款和待办共同复制，备份回退另走恢复准入\\
\bottomrule
\end{tabularx}
\caption{设计条件落实到路径、保存范围和恢复能力}\label{tab:ch16-design}
\end{table}

表中的速率和时限是设计输入。200 毫秒的受理目标覆盖入口到业务权威的有限路径，用户额外身份验证、支付机构等待和发货处理各有自己的阶段。后台恢复目标还依赖积压期间的持续到达率。例如消费能力每秒 2500 件而新事件每秒 2000 件，净排空能力只有每秒 500 件，清理 120 万件积压约需 40 分钟。要实现 20 分钟目标，净恢复能力至少每秒 1000 件，或在恢复期间收紧新责任的准入。

故障隔离由调用路径和保存责任共同形成。一个逻辑分区失去多数时，该分区暂停首次受理和需要新领取决定的付款发送。已有请求若在失去多数以前离开工作者，平台可以继续完成它，恢复后从原 C 查询结果。其他拥有多数的分区保持自己的受理能力。目录暂时不可用时，持有稳定路由缓存的入口可以直接访问原组，由数据组核验本地激活资格。需要首次发现或迁移新归属的请求则等待控制面恢复。

支付平台不可用会使 C/T 积压，却允许发布器继续传播此前已经保存的成功 EV。消息集群不可用时，付款结果与 Outbox 一起保留，用户可以读到已付款、履约传播等待中的进度。履约数据库不可用时，消费者保留投递或读取位置，后续在原事件处继续。图~\ref{fig:ch16-dependencies} 中的横向分支使这些暂停各有落点。与此同时，积压会占用磁盘、内存与扫描资源。系统按分区和责任类型设置准入、连接池与恢复并发，才能使一个外部依赖的长时间故障不经共享资源耗尽而拖慢其他路径。

\CSixteenFigure{ch16-13-dependencies}{不同依赖暂停时，工作停在各自的持久责任处，共享容量决定隔离能否维持}{fig:ch16-dependencies}

\subsection{稳定分区和部署故障域}

设 b1、b2、b3 为稳定的逻辑分区，G1、G2、G3 为当前承载它们的副本组。b1 当前放在 G1，以后三个副本可以更换，也可以把整个 b1 迁到 G2。\code{P:b1:1} 与 \code{C:b1:1} 中的 b1 表示身份归属，G1 表示当前执行位置。编号随分区事实保存，位置变化通过归属记录表达。

每个业务组拥有三个副本，控制组另有三个副本。图~\ref{fig:ch16-deployment} 把每组的一份副本放到三个设计地域。每一行形成自己的复制多数，行与行之间各有业务日志。控制组保存分区归属、配置和迁移资格，正常购买直接进入所属业务组。组合购买访问多个库存组，是因为它确实争用了多个资源。

\CSixteenFigure{ch16-02-deployment}{业务组与控制组各自跨三个设计地域形成多数，行间故障与恢复分别处理}{fig:ch16-deployment}

一台主机上的多个进程有独立地址、日志与存储目录，能表达进程退出、TCP 断连和复制切换。物理地域还会引入电力、机房、链路、磁盘和共同运维的故障相关性。设计中的三个地域要求各组的多数在相应故障条件下真实存活，消息集群、平台访问出口和消费者数据库也要满足自己的放置条件。业务副本跨域，外部支付入口或履约数据库的跨域能力则由其自身部署提供。

迁移属于控制路径。源组冻结 b1 的完整可变状态，导出一致图像，目标导入以后再交接归属资格并激活。在源被冻结、目标尚未激活的期间，b1 可以短暂停顿，b2 和 b3 继续。prepared 库存责任会影响这条路径：当前方案显式阻止带未终结准备责任的分区迁移，直到协调恢复已经完成。第十七章将把冻结范围和资格交接逐步展开。

\subsection{消息系统在全局架构中的位置}

购买权威形成付款结果以后，履约服务可能暂时不可用，处理速度也可能低于付款速度。消息系统在两者之间接收已付款事件，保存尚未完成的交付，使发布和消费按照各自的速度推进。发布器用确认判断消息域何时接过保存责任，消费者用确认或位移记录处理进度。第九、十章已经用这些接口连接过 Outbox 与 Inbox，本章从完整部署再看其内部结构。

RabbitMQ 的队列路径先由 exchange 按绑定路由到 queue，再把队列中尚未确认的消息交给消费者。消费者完成一条业务事务后确认该次投递，队列可据此释放相应交付责任。Kafka 的日志路径把事件追加到 topic 的某个分区，保留按 offset 定位的记录，各消费组分别保存读取位置。一个组消费完以后，日志还可以由另一个组独立读取，保留期限由日志政策决定~\cite{rabbitmq-queues,kafka-design}。

图~\ref{fig:ch16-message-overview} 展示两条路径的共同位置与不同进度模型。一次部署选用其中一个适配器，购买权威都提供同一份完整 EV，消费库都按 I/B 仲裁业务。队列投递标记和日志位置帮助高效传输与恢复，事件身份则随业务保存。选择时要结合订阅数量、历史重读、路由方式、积压规模与运维能力；后面的两节分别展开它们的副本、确认、并发和事务结构。

\CSixteenFigure{ch16-14-message-overview}{同一事件接口可接队列或分区日志，消费进度与业务身份各有保存位置}{fig:ch16-message-overview}

\subsection{沿一份业务事实贯穿全系统}

图~\ref{fig:ch16-facts} 把服务图换成事实关系图。Alice 的 $S$ 首先取得 \code{P:b1:1/C:b1:1/T:b1:1}，平台调用使用固定范围 \code{simulator:merchant-7:test:charge.v1} 与键 \code{pay-C:b1:1}。X1 是平台在自己的账本中保存的收款事实。购买权威核验这份事实以后保存 \code{O:b1:1} 和 \code{EV:b1:1}，消费者再把该订单的初次履约映射到 F501。

\CSixteenFigure{ch16-03-facts}{原身份、原付款事实、原事件与原履约通过持久关联连接}{fig:ch16-facts}

这条链可以有许多通信尝试。一份 EV 可能先占据 Kafka 分区中的 offset 0，发布器进程恢复后又占据 offset 1。RabbitMQ 也可能在确认响应丢失以后接收第二份同内容消息。它们继续引用同一个 EV。履约库按 EV 和初次履约业务守卫核验完整输入，新的投递位置继续指向原发货授权。

查询投影和缓存从这些权威事实派生。路由缓存记录 b1 当前在哪个组，进度缓存记录订单已付款或履约已接受，它们都可以重建。大图像存储保存分区状态的不可变块，以及描述完整块集合和应用位置的清单 manifest。恢复者根据完整清单取回图像，再核验它对应的权威前缀。第十八章讨论缓存的节点与集群结构，第十九章讨论大图像和事务存储，第二十章把可替换的执行者展开为批处理、任务、Actor 和函数。
